\documentclass[11pt]{article}
\usepackage[utf8]{inputenc}
\usepackage[T1]{fontenc}
\usepackage{amsmath}
\usepackage{amssymb}
\usepackage{multicol}
\usepackage{geometry}
\usepackage{tikz}
\usepackage{enumitem}
\usepackage{xcolor}
\usepackage{titlesec}

% Configurações de layout
\geometry{a4paper, left=1cm, right=1cm, top=0.5cm, bottom=1.2cm}
\setlength{\columnseprule}{0.4pt}

% Cores para títulos
\titleformat{\section}{\normalfont\Large\bfseries\color{blue}}{\thesection}{1em}{}
\titleformat{\subsection}{\normalfont\large\bfseries\color{red}}{\thesubsection}{1em}{}

\title{\textcolor{red}{Primeiro Trimestre: Combinatória
    \\Princípio Fundamental da Contagem, Arranjo, Permutação e Combinação}}
\author{Professor: Jefferson}
\date{}

\begin{document}

\maketitle
\vspace{-1cm}

\begin{center}
\large{Nome: \underline{\hspace{8cm}} \quad Turma: \underline{\hspace{3cm}}}
\end{center}

\begin{multicols}{2}

\section*{Exercícios - Combinatória}

\begin{enumerate}[leftmargin=*]

\item Uma loja vende camisetas em 5 cores diferentes e calças em 3 cores diferentes. Quantas combinações diferentes de camiseta e calça um cliente pode escolher?

\item Um restaurante oferece 4 opções de entrada, 6 opções de prato principal e 3 opções de sobremesa. De quantas maneiras diferentes um cliente pode escolher uma entrada, um prato principal e uma sobremesa?

\item Quantas placas de carro podem ser formadas com 3 letras (sem repetição) seguidas de 4 números (com repetição permitida)? Considere 26 letras no alfabeto e 10 algarismos.

\item De quantas maneiras 5 pessoas podem se organizar em uma fila?

\item Calcule o número de permutações de 8 objetos tomados 3 a 3, ou seja, $P(8,3)$.

\item De quantas formas diferentes 6 livros podem ser arrumados em uma prateleira?

\item Um código de segurança usa 4 dígitos diferentes. Quantos códigos diferentes podem ser formados?

\item Quantas maneiras diferentes existem para organizar as letras da palavra AMOR?

\item De quantas formas diferentes podemos escolher 3 alunos de uma turma de 10 alunos para uma comissão?

\item Calcule $C(10,3)$ (combinação de 10 elementos tomados 3 a 3).

\item De um grupo de 8 pessoas, quantas comissões de 4 pessoas podem ser formadas?

\item Quantas maneiras diferentes há para escolher 2 frutas de um conjunto de 7 frutas diferentes?

\item Em uma corrida com 10 participantes, de quantas formas diferentes podemos escolher o 1º, 2º e 3º lugar?

\item Calcule $A(7,2)$ (arranjo de 7 elementos tomados 2 a 2).

\item De quantas maneiras diferentes 5 atletas podem ocupar os 3 primeiros lugares de uma prova?

\item Uma pizzaria oferece 12 ingredientes diferentes. Se um cliente pode escolher até 3 ingredientes para sua pizza, quantas combinações diferentes ele pode fazer?

\item Quantas senhas de 6 caracteres podem ser formadas usando as 26 letras do alfabeto (com repetição permitida)?

\item De um baralho de 52 cartas, de quantas maneiras diferentes podemos escolher 5 cartas?

\item Em uma caixa há 4 bolas brancas, 3 bolas pretas e 2 bolas vermelhas. De quantas maneiras diferentes podemos escolher 3 bolas (sem considerar a ordem)?

\item Um sorteio premiará 2 pessoas entre 15 participantes. De quantas formas diferentes o sorteio pode ocorrer?

\item Quantas palavras de 4 letras (com sentido ou não) podem ser formadas com as letras A, B, C, D, E, F, sabendo que as letras não podem se repetir?

\item De um grupo de 9 pessoas, de quantas formas diferentes podemos escolher um presidente, um vice-presidente e um secretário?

\item A senha de um banco possui 5 dígitos. Quantas senhas diferentes podem ser criadas se os dígitos podem se repetir?

\item De quantas maneiras um professor pode escolher 4 alunos de uma classe de 30 para ajudantes?

\item Quantas maneiras diferentes há para organizar 6 pessoas em 3 pares, se a ordem dos pares não importa?

\item Um supermercado tem 15 tipos de sucos. De quantas formas diferentes um cliente pode escolher 3 tipos de sucos para comprar?

\item De quantas maneiras 4 pessoas podem se sentar ao redor de uma mesa redonda?

\item Calcule o número de permutações circulares de 6 objetos.

\item Uma delegação de 5 pessoas deve ser escolhida de um grupo de 20 pessoas. Quantas delegações diferentes podem ser formadas?

\item De um conjunto de 8 números, quantas sequências diferentes de 3 números (sem repetição) podem ser formadas?

\item Uma loja tem 10 estilos de camiseta, 5 estilos de calça e 3 estilos de sapato. Quantas combinações diferentes de roupa (1 camiseta, 1 calça e 1 sapato) um cliente pode escolher?

\item De quantas formas diferentes 7 CDs podem ser selecionados de uma coleção de 12 CDs?

\item Quantas maneiras diferentes há para distribuir 4 presentes diferentes entre 4 pessoas (cada pessoa recebe 1 presente)?

\item Um comitê de 5 pessoas será escolhido de um grupo contendo 6 mulheres e 4 homens. Quantos comitês diferentes podem ser formados se deve haver pelo menos 2 mulheres?

\item De um grupo de 12 pessoas, de quantas maneiras diferentes podemos dividir o grupo em 2 equipes de 6 pessoas cada (as equipes não têm rótulos)?

\item Quantas diagonais possui um polígono de 8 lados? (Dica: uma diagonal conecta dois vértices não adjacentes)

\end{enumerate}

\end{multicols}

\end{document}
